% Intended LaTeX compiler: pdflatex
\documentclass[a4paper,12pt]{article}
\usepackage[utf8]{inputenc}
\usepackage[T1]{fontenc}
\usepackage{graphicx}
\usepackage{longtable}
\usepackage{wrapfig}
\usepackage{rotating}
\usepackage[normalem]{ulem}
\usepackage{amsmath}
\usepackage{amssymb}
\usepackage{capt-of}
\usepackage{hyperref}
\usepackage{\string~/.emacs.d/common}
\author{mahmood}
\date{\today}
\title{premise}
\hypersetup{
 pdfauthor={mahmood},
 pdftitle={premise},
 pdfkeywords={},
 pdfsubject={},
 pdfcreator={Emacs 28.2 (Org mode 9.5.5)}, 
 pdflang={English}}
\begin{document}

\maketitle
\begin{note}
this document and others can be found on my website\\
\url{https://mahmoodsheikh36.github.io/}
\end{note}
\begin{definition}
a \textbf{premise} is a \href{20230401175139-statement.org}{proposition} on the basis of which we would be able to draw a \href{20230524212250-conclusion.org}{conclusion}, we can think of a premise as an evidence or assumption.\\
although n\\
\end{definition}
\begin{characteristic}
generally a premise is assumed to be true (although it doesnt have to be a tautology, as in it doesnt have to be universally true), as an argument is meaningful for its conclusion only when all of its premises are true. if one or more premises are false, the argument says nothing about whether the conclusion is true or false.\\
\begin{my_example}
"if x is an even integer then x\textsuperscript{2} is even" is a true statement, but x can also be an odd integer and the premise would be false.\\
you can make statements that are impossible, and prove them.\\
"If x is an integer that is both even and odd, then x\textsuperscript{2} is both even and odd"\\
The statement is true, and provable even though x can't exist.\\
\end{my_example}
\end{characteristic}
\end{document}